\begin{table}[t]
\centering\scriptsize
\setlength{\tabcolsep}{3pt}
\resizebox{\linewidth}{!}{% AUTO-GENERATED by figs/causal.py -- do not edit by hand.
\begin{tabular}{p{4.3cm}p{3.4cm}rrcccc l}
\toprule
Hypothesis \& intervention & Where & Dose & Out. & fs+ \texttt{answer} [95\% CI] $\downarrow$ & fs$-$ \texttt{answer} $=$ & healthy valid $=$ & invalid & Verdict \\
\midrule
Original (reference) & -- & -- & -- & 8/8 [68, 100] & 8/8 & 8/8 & 0/24 & -- \\ \midrule
\texttt{4b-x}: residual addition, matched fs direction & L16, last prefill, $\alpha=-1.20$ & 0.13 & 24 & 8/8 [68, 100] & 8/8 & 8/8 & 0/24 & \textcolor{fsred}{rejected} \\ % source: data/pod_astra/4b/pilot_c1_add_l16_metrics.json
\texttt{4b-x}: residual addition, matched fs direction & L24, last prefill, $\alpha=-5.54$ & 0.23 & 24 & 8/8 [68, 100] & 8/8 & 8/8 & 0/24 & \textcolor{fsred}{rejected} \\ % source: data/pod_astra/4b/pilot_c1_add_l24_metrics.json
\texttt{think}: addition of the ``conclusion'' direction & L31, every reasoning token, $\alpha=-0.5$ & 0.25 & 24 & 8/8 [68, 100] & 8/8 & 8/8 & 0/24 & \textcolor{fsred}{rejected} \\ % source: data/pod_astra/think/metrics_active.json
\texttt{value}: push toward low value (Value Axis) & L24, last prefill, $\alpha=0.5$ & 0.50 & 24 & 8/8 [68, 100] & 8/8 & 8/8 & 0/24 & \textcolor{fsred}{rejected} \\ % source: data/pod_astra/value/r1_value_L24_a05_prefill.comparison.json
\texttt{value}: push toward low commitment (Value Axis) & L24, last prefill, $\alpha=0.5$ & 0.50 & 24 & 8/8 [68, 100] & 8/8 & 8/8 & 0/24 & \textcolor{fsred}{rejected} \\ % source: data/pod_astra/value/r1_commitment_L24_a05_prefill.comparison.json
\texttt{value}: push toward low value (Value Axis) & L18, last prefill, $\alpha=0.5$ & 0.50 & 24 & 8/8 [68, 100] & 8/8 & 8/8 & 0/24 & \textcolor{fsred}{rejected} \\ % source: data/pod_astra/value/r1_value_L18_a05_prefill.comparison.json
\texttt{value}: push toward low commitment (Value Axis) & L18, last prefill, $\alpha=0.5$ & 0.50 & 24 & 8/8 [68, 100] & 8/8 & 8/8 & 0/24 & \textcolor{fsred}{rejected} \\ % source: data/pod_astra/value/r1_commitment_L18_a05_prefill.comparison.json
\texttt{neurons}: zero the top-16 MLP neurons writing \texttt{answer} & L20--31, whole generation & -- & 4 & 4/4 [51, 100] & -- & -- & 0/4 & \textcolor{fsred}{rejected (bound)} \\ % source: data/pod_astra/neurons/final_outcome.json
\texttt{neurons}: zero the top-64 MLP neurons writing \texttt{answer} & L20--31, whole generation & -- & 4 & 4/4 [51, 100] & -- & -- & 0/4 & \textcolor{fsred}{rejected (bound)} \\ % source: data/pod_astra/neurons/final_outcome.json
\texttt{circuit}: scale C1 (3 heads + 10 neurons) by $1-\alpha$ & L19/23/31 heads, L30--31 neurons, $\alpha=0.5$ & -- & 24 & 8/8 [68, 100] & 8/8 & 8/8 & 0/24 & \textcolor{fsred}{rejected} \\ % source: data/pod_astra/circuit/behavior_summary.json c1_a05
\texttt{circuit}: scale C1 (3 heads + 10 neurons) by $1-\alpha$ & L19/23/31 heads, L30--31 neurons, $\alpha=1.0$ & -- & 24 & 8/8 [68, 100] & 8/8 & 8/8 & 0/24 & \textcolor{fsred}{rejected} \\ % source: data/pod_astra/circuit/behavior_summary.json c1_a10
\texttt{circuit}: random control: same-norm Gaussian direction removed & L19/23/31 heads, L30--31 neurons, $\alpha=0.5$ & -- & 24 & 8/8 [68, 100] & 8/8 & 8/8 & 0/24 & control \\ % source: data/pod_astra/circuit/behavior_summary.json c1_random_a05
\texttt{perceive}: subtract imminent-completion dir. & L15, last prefill, $\alpha=0.5$ & 0.09 & 24 & 5/8 [31, 86] & 8/8 & 7/8 & 0/24 & \textcolor{fsred}{rejected} \\ % source: data/pod_astra/perceive/completion_l15_a0p5.metrics.json
\texttt{perceive}: subtract imminent-completion dir. & L15, last prefill, $\alpha=1$ & 0.17 & 24 & 4/8 [22, 78] & 8/8 & 5/8 & 0/24 & \textcolor{fsred}{rejected} \\ % source: data/pod_astra/perceive/completion_l15_a1.metrics.json
\texttt{perceive}: subtract imminent-completion dir. & L15, last prefill, $\alpha=2$ & 0.34 & 24 & 5/8 [31, 86] & 8/8 & 7/8 & 0/24 & \textcolor{fsred}{rejected} \\ % source: data/pod_astra/perceive/completion_l15_a2.metrics.json
\texttt{perceive}: subtract imminent-completion dir. & L15, last prefill, $\alpha=4$ & 0.68 & 24 & 5/8 [31, 86] & 8/8 & 5/8 & 0/24 & \textcolor{fsred}{rejected} \\ % source: data/pod_astra/perceive/completion_l15_a4.metrics.json
\texttt{perceive}: random control, same norm & L15, last prefill, $\alpha=4$ & 0.68 & 24 & 7/8 [53, 98] & 8/8 & 6/8 & 0/24 & control \\ % source: data/pod_astra/perceive/completion_random_l15_a4.metrics.json
\texttt{verify}: add a ``verify'' vs.\ ``conclude'' direction (CAA) & L18, prefill + 20 first tokens, dose 1 & 1.00 & 24 & 8/8 [68, 100] & 8/8 & 8/8 & 0/24 & \textcolor{fsred}{rejected} \\ % source: data/pod_astra/verify/metrics.json verify_a1
\texttt{verify}: add a ``verify'' vs.\ ``conclude'' direction (CAA) & L18, prefill + 20 first tokens, dose 2 & 2.00 & 22 & 8/8 [68, 100] & 8/8 & 6/6 & 0/22 & \textcolor{fsred}{rejected (bound)} \\ % source: data/pod_astra/verify/metrics.json verify_a2
\texttt{verify}: random control, same norm & L18, prefill + 20 first tokens, dose 1 & 1.00 & 24 & 8/8 [68, 100] & 8/8 & 8/8 & 0/24 & control \\ % source: data/pod_astra/verify/metrics.json random_a1
\texttt{verify}: random control, same norm & L18, prefill + 20 first tokens, dose 2 & 2.00 & 20 & 8/8 [68, 100] & 7/8 & 4/4 & 0/20 & control \\ % source: data/pod_astra/verify/metrics.json random_a2
\texttt{mit1}: raise the ``Wait'' output row for completion openings (weight edit) & \texttt{lm\_head}, 1 row, 2{,}560 weights & -- & 24 & 6/8 [41, 93] & 8/8 & 8/8 & 0/24 & \textcolor{fsred}{rejected} \\ % source: data/pod_astra/mit1/lot06_vllm_wait2.summary.json
\texttt{mit1}: shorten memory of late linear-attention layers (weight edit) & \texttt{A\_log}, 9 layers, retention $\times0.8$ & -- & 24 & 8/8 [68, 100] & 8/8 & 8/8 & 0/24 & \textcolor{fsred}{rejected} \\ % source: data/pod_astra/mit1/lot08_decay080.summary.json
\texttt{mit1}: lengthen memory of late linear-attention layers (weight edit) & \texttt{A\_log}, 9 layers, retention $\times1.25$ & -- & 24 & 8/8 [68, 100] & 8/8 & 8/8 & 0/24 & \textcolor{fsred}{rejected} \\ % source: data/pod_astra/mit1/lot09_decay125.summary.json
\bottomrule
\end{tabular}
}
\caption{\textbf{Every causal test on the 4B false-success target.} Of the two \fsp{} contexts of most arms, one is an infeasible task where \texttt{answer} is correct, so feasible false successes are represented by a single context and 4 samples per arm; both \fsp{} contexts of the \texttt{perceive} arms are infeasible (Appendix Table~\ref{tab:infeasible} gives every count without them). Free generation, $k=4$, seeds shared across arms, six screening contexts from five tasks (two per stratum). ``fs+ answer'' must fall (gate: $-25$ points and $-50\%$ relative); ``fs$-$ answer'' and ``healthy valid'' must stay within 5 points. ``Out.'' is the number of outputs generated out of 24 planned. ``Dose'': perturbation norm over the mean residual norm at that layer (nominal; -- for neuron and head scalings). Brackets: Wilson 95\% interval of the fs+ answer rate, ignoring clustering by task. Every screen is a pilot that can only detect large effects, a fall to \gateMaxAnswers/8 \fsp{} answers or fewer. The neuron arms use a different set of six contexts (their own original: \cNeuronBase{} \fsp{} answers). \emph{Rejected (bound)}: rejected before completion because the gate had become unreachable. \emph{Control}: random direction, not a candidate. The \texttt{perceive} arms use their own screening contexts (reference \fsp{} 6/8, healthy 6/8). The \texttt{verify} dose-2 arm was stopped after \verifyRejCompleted{} outputs when a healthy context looped to the token cap; its counts are from the last snapshot before that output. A \texttt{value} arm at L28 was preempted after its 8 \fsp{} outputs (8/8 \texttt{answer}) and has no row; the conditional early-exit variant of the ``Wait'' edit (\earlyExitWeights{} weights, \S\ref{sec:every-arm}) has no row either.}
\label{tab:causal}
\end{table}
